\documentclass[12pt, reqno]{amsart}
\usepackage{latexsym, amssymb, amscd, xypic, amsthm, setspace}
\input xy
\xyoption{all}

%%: Font = times
\usepackage{mathptmx} % rm & math
\usepackage[scaled=0.90]{helvet} % ss
\usepackage{courier} % tt
\normalfont
\usepackage[T1]{fontenc}
\pagestyle{plain}


\newtheorem{theorem}{Theorem}[subsection]
\newtheorem{prop}[theorem]{Proposition}
\newtheorem{lemma}[theorem]{Lemma}
\newtheorem{corollary}[theorem]{Corollary}
\newtheorem{claim}[theorem]{Claim}

\theoremstyle{definition}
\newtheorem{de}[theorem]{Definition}
\newtheorem{notation}{Notation}
\newtheorem{example}[theorem]{Example}
\newtheorem{examples}{Examples}
\newtheorem{remark}[theorem]{Remark}
\newtheorem{remarks}{Remarks}
\newtheorem{question}{Question}


\newtheorem*{nthm}{Theorem}
\newtheorem*{nlem}{Lemma}
\newtheorem*{nprop}{Proposition}
\newtheorem*{ncor}{Corollary}
\newtheorem*{nclaim}{Claim}
\newtheorem*{nde}{Definition}
\newtheorem*{nex}{Example}
\newenvironment{solution}{\par\smallskip\noindent\textbf{Solution.}}{\hfill$\square$\par\medskip}

\newtheoremstyle{test}{3pt}{3pt}{\itshape}{\parindent}{\bfseries}{---}{  }{\thmnote{#3}} \theoremstyle{test}
\newtheorem*{mythm}{}


%: paragraphs

\newcommand{\point}{\vspace{3mm}\par \noindent (\refstepcounter{subsection}{\bf \thesubsection) } }

\newcommand{\tpoint}[1]{\vspace{3mm}\par \noindent \refstepcounter{subsubsection}{\bf \thesubsubsection.}
  {\em #1. ---} }
\newcommand{\epoint}[1]{\vspace{3mm}\par \noindent \refstepcounter{subsection}{\bf \thesubsection.}
  {\em #1.} }
\newcommand{\spoint}{\vspace{3mm}\par \noindent \refstepcounter{subsection}{\bf \thesubsection.} }

\newcommand{\bpoint}[1]{\vspace{3mm}\par \noindent \refstepcounter{subsection}{\bf \thesubsection.}
  {\bf #1.} }



%: numbering and margins

\numberwithin{equation}{subsection}

\setlength{\oddsidemargin}{0cm} \setlength{\evensidemargin}{0cm}
\setlength{\marginparwidth}{0in}
\setlength{\marginparsep}{0in}
\setlength{\marginparpush}{0in}
\setlength{\topmargin}{0in}
\setlength{\headheight}{0pt}
\setlength{\headsep}{0pt}
\setlength{\footskip}{.3in}
\setlength{\textheight}{9.2in}
\setlength{\textwidth}{6.5in}
\setlength{\parskip}{0.25pt}
\setlength{\parindent}{0.25in}

%: math shortcuts


\DeclareMathOperator{\inv}{inv}
\DeclareMathOperator{\out}{Out}
\DeclareMathOperator{\gal}{Gal}
\renewcommand{\mod}{\mathrm{\, mod \, }}
\renewcommand{\d}[1]{{}^{#1}}
\renewcommand{\b}[1]{\mathbf{{#1}}}
\newcommand{\C}{\mathbb{C}}
\newcommand{\be}[1]{\begin{eqnarray} \label{#1}}
\newcommand{\ee}{\end{eqnarray}}
\newcommand{\re}{\mathsf{Re}}
\newcommand{\im}{\mathsf{Im}}
\newcommand{\rr}{\rightarrow}

\newcommand{\n}{\mathbb{N}}
\newcommand{\zee}{\mathbb{Z}}
\newcommand{\R}{\mathbb{R}}
\newcommand{\pr}{\mathbb{P}}
\renewcommand{\gcd}[2]{\mathsf{gcd}(#1, #2)}
\newcommand{\ov}[1]{\overline{#1}}
\renewcommand{\Re}{\mathrm{Re}}
\renewcommand{\Im}{\mathrm{Im}}



\title{Problem Set 1: Solutions \\ MATH 411: Honors Complex Variables, Fall 2026}
\date{last compiled: \today}

\numberwithin{equation}{section}
\refstepcounter{section}

\begin{document}						%%
\maketitle



\begin{enumerate}

\item
\begin{enumerate}

\item Put $-1 - i$ into polar form
\begin{solution}
$|-1-i|=\sqrt2$, and an argument is $5\pi/4$ (or $-3\pi/4$).  Hence
\[ -1-i=\sqrt2e^{5\pi i/4}=\sqrt2e^{-3\pi i/4}. \]
\end{solution}
\item Put $e^{- 5 i \pi /4}$ into the form $x + i y$.
\begin{solution}
By Euler's formula,
\[e^{-5\pi i/4}=\cos(-5\pi/4)+i\sin(-5\pi/4)
=-\frac1{\sqrt2}+\frac{i}{\sqrt2}.\]
\end{solution}
\item Find the real and imaginary parts of $i^{1/4}$, taking the fourth root whose angle lies between $0$ and $\pi/2$.
\begin{solution}
The required root is $e^{i\pi/8}$, so
\[i^{1/4}=\frac{\sqrt{2+\sqrt2}}2+i\frac{\sqrt{2-\sqrt2}}2.\]
Thus its real and imaginary parts are $\sqrt{2+\sqrt2}/2$ and
$\sqrt{2-\sqrt2}/2$, respectively.
\end{solution}
\item Describe geometrically the set of points $z \in \mathbb{C}$ such that $|z| = \Re(z)+1$.
\begin{solution}
Writing $z=x+iy$ and squaring gives
\[x^2+y^2=(x+1)^2\quad\Longleftrightarrow\quad y^2=2x+1.\]
This is the parabola $x=(y^2-1)/2$, with vertex $(-1/2,0)$ and opening
to the right.  Its points have $x+1\geq1/2$, so squaring added no points.
\end{solution}
\end{enumerate}
\item Let $f(z) = e^z$. Describe the image under $f$ of the following sets
\begin{enumerate}
\item The set of $z = x + i y $ such that $x \leq 1$ and $0 \leq y \leq \pi$
\begin{solution}
Since $e^{x+iy}=e^xe^{iy}$, the modulus ranges over $(0,e]$ and the
argument over $[0,\pi]$.  The image is
\[\{\zeta\in\mathbb C:0<|\zeta|\leq e,\ \Im\zeta\geq0\},\]
the closed upper half-disc of radius $e$, with the origin removed.
\end{solution}
\item The set of $z = x + i y $ such that $0 \leq y \leq \pi$ (no condition on $x$)
\begin{solution}
Here the modulus ranges over all positive numbers, so the image is
\[\{\zeta\in\mathbb C:\zeta\neq0,\ \Im\zeta\geq0\},\]
the closed upper half-plane with the origin removed.
\end{solution}
\end{enumerate}

\item Let $f$ be the function defined by $$ f(z) = \lim_{ n \rr \infty} \frac{1}{1 + n^2 z} .$$
 Show that $f$ is the characteristic function of the set $\{ 0 \}$, i.e. $f(0)=1$ and $f(z)=0$ otherwise.
\begin{solution}
For $z=0$, every term is $1$.  If $z\neq0$, then, once $n^2|z|>1$,
\[\left|\frac1{1+n^2z}\right|\leq\frac1{n^2|z|-1}\longrightarrow0.\]
Therefore $f(0)=1$ and $f(z)=0$ for every $z\neq0$.
\end{solution}

 \item \begin{enumerate}
\item Let $z, w$ be two complex numbers such that $\overline{z} w \neq 1$. Prove that $$ \left| \frac{w -z }{1 - \overline{w} z} \right| < 1 \text{ if } |z| < 1 \text{ and } |w| < 1 $$ and also that $$ \left| \frac{w -z }{1 - \overline{w} z} \right| = 1 \text{ if } |z| = 1 \text{ or } |w| = 1 $$

\emph{Hint: reduce to the case where $z$ is real, say $z=r$, in which case it suffices to prove that $$ (r - w) (r - \overline{w}) \leq ( 1- rw) (1 - r \overline{w}) $$ with equality when $|r|=1$ or $|w|=1$.}

\begin{solution}
Expanding both sides gives
\[
|1-\overline wz|^2-|w-z|^2=(1-|w|^2)(1-|z|^2).
\]
If $|z|<1$ and $|w|<1$, the right side is positive, and division by
$|1-\overline wz|^2$ gives the strict inequality.  If $|z|=1$ or
$|w|=1$, the right side is zero, so the quotient has modulus $1$.
\end{solution}

\item Prove that for a fixed $w$ in the unit disc $\mathbb{D}$ the mapping $$ F: z \mapsto \frac{w - z}{1 - \overline{w} z} $$ satisfies the following conditions:

\begin{enumerate}
\item $F$ maps $\mathbb{D}$ to itself and is holomorphic on $\mathbb{D}$
\item $F$ interchanges $0$ and $w$ (i.e. $F(0)=w$ and $F(w)=0$)
\item $|F(z)|=1$ if $|z|=1$
\item $F: \mathbb{D} \rr \mathbb{D}$ is bijective (Hint: calculate $F \circ F$).
\end{enumerate}

\begin{solution}
Since $|\overline wz|<1$ for $z,w\in\mathbb D$, the denominator never
vanishes in $\mathbb D$; hence $F$ is holomorphic there.  The preceding
part shows that $|F(z)|<1$, so $F(\mathbb D)\subseteq\mathbb D$.  Also
\[F(0)=w,\qquad F(w)=0,\]
and the preceding part gives $|F(z)|=1$ when $|z|=1$.

Finally, direct substitution yields
\[
F(F(z))=
\frac{w-\dfrac{w-z}{1-\overline wz}}
{1-\overline w\dfrac{w-z}{1-\overline wz}}=z.
\]
Thus $F$ is its own inverse and is therefore a bijection of $\mathbb D$.
\end{solution}

\end{enumerate}

\item Show that in polar coordinates, i.e. if $F(x, y)$ is regarded as $F(r, \theta)= ( u(r,\theta), v(r, \theta))$, the Cauchy--Riemann equations take the form  $$ \frac{\partial u}{\partial r}  = \frac{1}{r} \frac{ \partial v}{\partial \theta} \text{ and } \frac{1}{r} \frac{ \partial u}{\partial \theta} = - \frac{\partial v}{\partial r}.$$  Use these equations to show that the logarithm function defined as $$ \log(z) := \log (r) + i \theta \text{ where } z = r e^{ i \theta} \text{ with } - \pi < \theta < \pi $$ is holomorphic in the region $r > 0$ and $- \pi < \theta < \pi$.
\begin{solution}
The chain rule gives
\[
u_r=u_x\cos\theta+u_y\sin\theta,\qquad
u_\theta=-ru_x\sin\theta+ru_y\cos\theta,
\]
and the analogous formulas for $v$.  Substituting $u_x=v_y$ and
$u_y=-v_x$ gives
\[u_r=\frac1r v_\theta,\qquad \frac1r u_\theta=-v_r.\]

For the given logarithm, $u=\log r$ and $v=\theta$.  Thus
$u_r=1/r$, $u_\theta=0$, $v_r=0$, and $v_\theta=1$, so both polar
Cauchy--Riemann equations hold.  The partial derivatives are continuous
for $r>0$ and $-\pi<\theta<\pi$, so this branch is holomorphic there.
\end{solution}

\item \begin{enumerate}
\item Show that $$4 \frac{\partial}{\partial z} \frac{\partial}{\partial \overline{z}} = 4 \frac{\partial}{\partial \overline{z}} \frac{\partial}{\partial z} = \Delta $$ where $\Delta$ is the Laplacian, $$ \Delta = \frac{\partial^2}{\partial x^2} + \frac{\partial^2}{\partial y^2} .$$
\begin{solution}
Using
\[\partial_z=\frac12(\partial_x-i\partial_y),\qquad
\partial_{\overline z}=\frac12(\partial_x+i\partial_y),\]
and the equality of mixed partials, we obtain
\[
4\partial_z\partial_{\overline z}
=(\partial_x-i\partial_y)(\partial_x+i\partial_y)
=\partial_x^2+\partial_y^2=\Delta.
\]
Reversing the order gives the same calculation.
\end{solution}
\item Show that if $f$ is holomorphic on an open set $U \subset \mathbb{C}$, then $\psi(x, y) = \Re(f)$ and $\phi(x, y) = \Im(f)$ are both harmonic functions (i.e. they satisfy $\Delta \psi = \Delta \phi = 0$).
\begin{solution}
The Cauchy--Riemann equations are $\psi_x=\phi_y$ and
$\psi_y=-\phi_x$.  Since holomorphic functions are smooth,
\[
\Delta\psi=\psi_{xx}+\psi_{yy}=\phi_{yx}-\phi_{xy}=0,
\qquad
\Delta\phi=\phi_{xx}+\phi_{yy}=-\psi_{yx}+\psi_{xy}=0.
\]
Thus both real and imaginary parts are harmonic.
\end{solution}
\end{enumerate}

\item Show that $f(x+ i y) = \sqrt{|x| | y |}$ satisfies the Cauchy-Riemann equations at the origin but that $f$ is not holomorphic at $0$
\begin{solution}
Write $f=u+iv$, where $u(x,y)=\sqrt{|x||y|}$ and $v=0$.  On either
coordinate axis $u$ is zero, so
\[u_x(0,0)=u_y(0,0)=v_x(0,0)=v_y(0,0)=0.\]
Thus the Cauchy--Riemann equations hold at the origin.  But along the real
axis the difference quotient is zero, whereas for $t>0$,
\[
\frac{f(t+it)-f(0)}{t+it}=\frac{t}{t(1+i)}=\frac1{1+i}.
\]
The limits disagree, so $f$ is not complex differentiable at $0$ and
hence is not holomorphic there.
\end{solution}

\item Suppose that $f$ is holomorphic in an open disc $\mathcal{D} \subset \mathbb{C}$, say $f= u(x, y) +  i v(x, y)$. Prove that \begin{enumerate}
\item if $|f|^2 = u^2 + v^2$ is constant, then $f$ is constant;
\begin{solution}
Suppose $|f|^2=c$. If $c=0$ we have nothing to show, so assume
$c=u^2+v^2>0$. Differentiating the expression $|f|^2=c$ yields
\begin{eqnarray*}
2uu_x+2vv_x&=&0,\\
2uu_y+2vv_y&=&0.
\end{eqnarray*}
Applying Cauchy--Riemann (i.e. $u_x=v_y$, $u_y=-v_x$), we find
\begin{eqnarray*}
2uu_x-2vu_y&=&0,\\
2uu_y+2vu_x&=&0.
\end{eqnarray*}
Since
\[
\det\begin{pmatrix}u&-v\\v&u\end{pmatrix}=u^2+v^2\neq0,
\]
the only solution to this system is $u_x=u_y=0$. In a similar way we
also conclude that $v_x=v_y=0$. As the Jacobian of the map
$F=(u(x,y),v(x,y))$ vanishes, the map $F$ is locally constant, and hence
constant since $\mathcal D$ is connected.
\end{solution}
\item if there exist real, non-zero constants $a, b,$ and $c$ such that $$ a u (x, y) + b v (x, y) = c, $$ then $f$ is constant.
\begin{solution}
Differentiating the above expression and applying Cauchy--Riemann yields
the system
\[
\begin{pmatrix}a&b\\b&-a\end{pmatrix}
\begin{pmatrix}u_x\\v_x\end{pmatrix}
=\begin{pmatrix}0\\0\end{pmatrix}.
\]
As
\[
\det\begin{pmatrix}a&b\\b&-a\end{pmatrix}=-(a^2+b^2)\neq0,
\]
we must have $u_x=v_x=0$. Again applying Cauchy--Riemann, we also
conclude that $u_y=v_y=0$, and we argue as above.
\end{solution}
\end{enumerate}

\end{enumerate}


\end{document}
